% !TEX program = lualatex
\documentclass[11pt,a4paper]{ltjsarticle}
\usepackage{icat-common}

% 追加パッケージ
\usepackage{bm}

% 定理環境
\theoremstyle{definition}
\newtheorem{theorem}{定理}[section]
\newtheorem{lemma}[theorem]{補題}
\newtheorem{proposition}[theorem]{命題}
\newtheorem{corollary}[theorem]{系}
\newtheorem{definition}[theorem]{定義}
\newtheorem{example}[theorem]{例}
\newtheorem{remark}[theorem]{注記}

% コマンド定義
\newcommand{\Md}{\mathbf{D}}
\newcommand{\Mnn}{\mathbf{NN}}
\newcommand{\Mr}{\mathbf{R}}
\newcommand{\MC}{\mathbf{C}}
\newcommand{\MT}{\mathbf{T}}
\newcommand{\phiY}{\phi_{\text{Y}}}
\newcommand{\phiN}{\phi_{\text{N}}}
\newcommand{\phiYN}{\phi_{\text{YN}}}
\newcommand{\phiNS}{\phi_{\text{NS}}}
\newcommand{\Herbrand}{\mathcal{H}}
\newcommand{\Herbase}{\mathcal{B}}
\newcommand{\Res}{\mathsf{Res}}
\newcommand{\LLM}{\mathsf{LLM}}
\newcommand{\Prob}{\mathbb{P}}

\title{モナド五蘊の洗練モナド（想）における\\エルブラン空間探索層の導入と\\論理的・確率論的探索の統一的扱い}
\author{千葉将臣}
\date{2026-07-03}

\begin{document}

\maketitle

\begin{abstract}
本論文は、顕現論のモナド五蘊（$\Md, \Mnn, \Mr, \MC, \MT$）のうち、洗練モナド（想）$\Mr$ にエルブラン空間（Herbrand universe）での探索という層を導入することで、論理的探索（導出・バックトラッキング）と確率論的探索（LLMのような確率モデル）を統一的に扱う枠組みを提案する。エルブラン空間を探索空間と見なし、導出戦略を探索の効率化手法として捉えることで、哲学（認識論）からプログラミング（アルゴリズム設計）までを一貫した枠組みで記述できることを示す。さらに、LLMとの接続は想蘊に限定されないことを確認し、$\Md, \Mnn, \Mr, \MC, \MT$ の各モナドに局所的な確率論的知性を差し込む発展的構想を整理する。また、この枠組みが顕現論の四句分別（是・非・亦是亦非・非是非非）およびモナド五蘊の再帰構造と整合的であることを論じる。
\end{abstract}

\section{はじめに}

顕現論（顕現論・界論・無記の意味論）は、モナド五蘊（$\Md, \Mnn, \Mr, \MC, \MT$）と四句分別（$\phiY, \phiN, \phiYN, \phiNS$）を基礎として、認識・計算・存在を統一的に扱う枠組みを目指している。本論文では、特に五蘊の「想」に対応する洗練モナド $\Mr$ に注目し、エルブラン空間（Herbrand universe）での探索という層を導入することで、論理的探索（導出・バックトラッキング）と確率論的探索（LLMのような確率モデル）を統一的に扱うことを提案する。そのうえで、LLMを接続する対象を $\Mr$ に限定せず、モナド五蘊全体に拡張する発展的可能性も検討する。

エルブラン空間は一階述語論理における項の集合であり、証明探索やモデル理論の基礎として用いられる\cite{herbrand}。本論文では、エルブラン空間を「探索空間」と見なし、導出戦略を「探索の効率化手法」として捉えることで、哲学（認識論）からプログラミング（アルゴリズム設計）までを一貫した枠組みで記述できることを示す。

\section{背景}

\subsection{モナド五蘊と四句分別}

顕現論では、五蘊（色・受・想・行・識）をモナドとして再解釈し、以下のように対応づける\cite{icat}。

\begin{itemize}
\item $\Md$（色）：搾取モナド（証明障害の残余）
\item $\Mnn$（受）：二重否定モナド（$\neg\neg$ の残余）
\item $\Mr$（想）：洗練モナド（双方向演繹）
\item $\MC$（行）：カット消去モナド（証明の正規化）
\item $\MT$（識）：統制モナド（識軸スパイラル $S_{\MT}$）
\end{itemize}

また、四句分別は以下のように定義される\cite{icat}。

\begin{itemize}
\item 是（$\phiY$）：対象そのもの
\item 非（$\phiN$）：否定された対象
\item 亦是亦非（$\phiYN$）：否定の双方向性
\item 非是非非（$\phiNS$）：不可能性の残余（決定不能性、ランダム性、外在対象の内部反映など）
\end{itemize}

本論文では、特に $\Mr$（想）を「世界をどのように構造化し、探索するかを決めるプロセス」として捉え直す。

\subsection{エルブラン空間と証明探索}

エルブラン空間 $\Herbrand$ は、一階述語論理の項の集合であり、エルブラン基底 $\Herbase$ は原子論理式の集合として定義される\cite{herbrand}。証明探索は、エルブラン基底上の節集合から矛盾を導き出すプロセスであり、導出（resolution）はその基本手法である\cite{resolution}。

導出戦略（SLD導出、線形導出、単位導出など）は、どの節をどの順番で適用するかを制御する戦略であり、探索の効率化手法として理解できる\cite{resolution}。

\subsection{LLMと確率論的探索}

LLM（大規模言語モデル）は、テキストの確率分布を学習し、次のトークンを確率的に生成するモデルである\cite{llm}。エルブラン空間の観点では、LLMは探索空間上の確率分布を学習し、どの分岐を優先的に探索するかを確率的に決定するモデルと見なせる。

\section{提案：エルブラン空間探索層の導入}

\subsection{洗練モナド $\Mr$ の再定義}

洗練モナド $\Mr$（想）を、エルブラン空間 $\Herbrand$ 上の探索プロセスとして定義する。

\begin{definition}[洗練モナド $\Mr$]
圏 $\mathcal{C}$ を基礎圏とし、エルブラン空間 $\Herbrand$ とエルブラン基底 $\Herbase$ を固定する。洗練モナド $\Mr$ は、以下の三つ組 $(T_{\Mr}, \eta_{\Mr}, \mu_{\Mr})$ として定義される。
\begin{itemize}
\item 関手 $T_{\Mr} : \mathcal{C} \to \mathcal{C}$：探索空間上の構造を反映する関手
\item 単位 $\eta_{\Mr} : \mathrm{Id} \Rightarrow T_{\Mr}$：探索の開始（初期状態）
\item 乗法 $\mu_{\Mr} : T_{\Mr}^2 \Rightarrow T_{\Mr}$：探索の繰り返し（再帰的探索）
\end{itemize}
\end{definition}

$\Mr$ は、エルブラン空間 $\Herbrand$ 上の探索戦略（導出戦略・確率論的戦略）を統一的に扱うモナドとして機能する。

\subsection{論理的探索と確率論的探索の統合}

エルブラン空間 $\Herbrand$ を探索空間と見なし、以下の二つの探索を統一的に扱う。

\subsubsection{論理的探索（導出・バックトラッキング）}

論理的探索は、ゴールから出発して深さ優先探索を行い、制約が失敗したらバックトラックするプロセスである。導出戦略（SLD導出など）は、探索の効率化手法として働く\cite{resolution}。

エルブラン空間の観点では、論理的探索は以下のように定義できる。

\begin{definition}[論理的探索]
エルブラン空間 $\Herbrand$ 上の論理的探索は、有向グラフ $G = (V, E)$ として定義される。
\begin{itemize}
\item 頂点 $V$：部分証明（ゴール集合）
\item 辺 $E$：推論規則の適用（導出ステップ）
\end{itemize}
探索戦略は、辺の選択順序を制御する関数 $s : V \to E^*$ として定義される。
\end{definition}

\subsubsection{確率論的探索（LLMのようなモデル）}

確率論的探索は、LLMのような確率モデルを用いて、どの分岐を優先的に探索するかを確率的に決定するプロセスである\cite{llm}。

エルブラン空間の観点では、確率論的探索は以下のように定義できる。

\begin{definition}[確率論的探索]
エルブラン空間 $\Herbrand$ 上の確率論的探索は、確率分布 $\Prob : V \to [0,1]$ を持つ有向グラフ $G = (V, E)$ として定義される。LLMは、この確率分布を学習し、次の頂点 $v'$ を $\Prob(v' \mid v)$ に従って選択する。
\end{definition}

\subsubsection{統合の枠組み}

洗練モナド $\Mr$ は、論理的探索と確率論的探索を統一的に扱うためのインターフェースとして機能する。

\begin{definition}[統合的探索モナド]
洗練モナド $\Mr$ は、以下のように定義される。
\[
T_{\Mr}(X) = \Res(X) \oplus \LLM(X)
\]
ここで、
\begin{itemize}
\item $\Res(X)$：論理的探索（導出・バックトラッキング）による探索空間
\item $\LLM(X)$：確率論的探索（LLM）による探索空間
\item $\oplus$：探索空間の統合（直和またはより高度な結合）
\end{itemize}
単位 $\eta_{\Mr}$ と乗法 $\mu_{\Mr}$ は、探索の開始と再帰的探索を定義する。
\end{definition}

この枠組みにより、例えば以下のようなハイブリッド探索が可能になる。

\begin{example}[ハイブリッド探索]
\begin{enumerate}
\item LLMが有望な分岐を確率的に提案する（$\LLM(X)$）。
\item 論理エンジンがその分岐に対して厳密な証明を行う（$\Res(X)$）。
\item 失敗した場合はバックトラックし、別の分岐を試す。
\end{enumerate}
\end{example}

\section{顕現論との整合性}

\subsection{四句分別との対応}

エルブラン空間探索層は、四句分別と以下のように対応する\cite{icat}。

\begin{itemize}
\item 是（$\phiY$）：エルブラン空間 $\Herbrand$（可能世界の全体）
\item 非（$\phiN$）：探索の失敗・バックトラック
\item 亦是亦非（$\phiYN$）：双方向の探索プロセス（論理的探索）
\item 非是非非（$\phiNS$）：確率的不確実性・残余（確率論的探索）
\end{itemize}

\subsection{モナド五蘊の再帰構造}

モナド五蘊の再帰構造（$\Md \to \Mnn \to \Mr \to \MC \to \MT \to \dots$）は、探索戦略の洗練・最適化・統制のプロセスとして理解できる。

\begin{itemize}
\item $\Mr$（想）：探索戦略の洗練（エルブラン空間探索層）
\item $\MC$（行）：探索の最適化（カット消去・枝刈り）
\item $\MT$（識）：探索の統制（識軸スパイラル $S_{\MT}$）
\end{itemize}

さらに、この再帰構造は単なる線形系列ではなく、
\[
\Md \to \Mnn \to \Mr \to \MC \to \MT \to \Md \to \dots
\]
という循環構造として理解できる。このため、$\Mr$ にLLMを接続することは、五蘊全体のうち想蘊だけに確率論的知性を与える特殊な場合であり、同じ発想は他のモナドにも拡張できる。

\section{発展：モナド五蘊全体へのLLM接続}

本論文の中心は、洗練モナド $\Mr$ におけるエルブラン空間探索層である。しかし、LLMとの接続は想蘊に限定される必要はない。モナド五蘊は再帰的かつ循環的に結合しているため、各モナドに局所的な確率論的知性を差し込み、それらを $\MT$ によって統制する分散型の構成が考えられる。

\subsection{$\Md$（色）：搾取モナドへの接続}

$\Md$ は、証明障害の残余、すなわち証明過程において回収しきれず、搾取として残る構造を扱うモナドである。ここにLLMを接続する場合、LLMは証明探索中にどの部分が障害になりやすいかを確率的に予測する補助機構として働く。

具体的には、LLMは過去の証明失敗、局所的な式変形、導出不能な節の形、探索木の急激な分岐などを手がかりとして、搾取を生みやすい構造をパターン認識する。このときLLMは証明そのものを置き換えるのではなく、証明障害が発生しやすい場所を先取りして示す確率論的センサーとして機能する。

\begin{example}[$\Md$ における障害予測]
証明探索中に、ある補題の適用後だけ導出木が急激に分岐し、閉じない枝が多数生じるとする。このときLLMは、その補題周辺に搾取の残余が蓄積している可能性を確率的に警告し、別の補題選択や局所的な定式化の修正を促す。
\end{example}

\subsection{$\Mnn$（受）：二重否定モナドへの接続}

$\Mnn$ は、$\neg\neg$ の残余、すなわち否定の否定として内部化される構造を扱うモナドである。ここにLLMを接続すると、LLMはどの否定が二重否定として回収可能か、またどの否定が四句分別における亦是亦非（$\phiYN$）として双方向的に扱えるかを確率的に探索する。

この接続では、LLMは否定の形式的変換を決定する主体ではない。むしろ、否定の配置、前提の依存関係、反証可能性の形を観察し、「この否定は二重否定として扱える可能性が高い」という候補を提示する。厳密な可否は、その後に論理エンジンまたは型検査的な検証過程によって判定される。

\begin{example}[$\Mnn$ における二重否定候補]
ある主張 $A$ に対して直接証明が難しい一方、$\neg A$ を仮定すると矛盾が導かれる見込みが高い場合、LLMは $\neg\neg A$ を中間目標として導入することを提案する。
\end{example}

\subsection{$\Mr$（想）：洗練モナドへの接続}

$\Mr$ へのLLM接続は、本論文の中心的構成である。ここでは、LLMはエルブラン空間上の分岐に確率的な優先順位を与え、導出戦略の動的選択を支援する。すなわち、LLMは「どの節を先に試すべきか」「どの項の具体化が有望か」「どの導出経路が閉じやすいか」を確率的に提案する。

この接続により、論理的探索は完全性や健全性を保持する検証層として残りつつ、探索順序の選択には確率論的な柔軟性を導入できる。LLMは証明を確定するのではなく、有望な分岐を提示し、論理エンジンがそれを厳密に検証する。

\subsection{$\MC$（行）：カット消去モナドへの接続}

$\MC$ は、証明の正規化、すなわちカット消去や枝刈りを通じて証明構造を簡約するモナドである。ここにLLMを接続する場合、LLMはどのカットを消去すべきか、どの中間補題が証明全体を複雑化しているか、どの枝が冗長であるかを確率的に推定する。

たとえば、LLMは複数の証明候補を比較し、同じ結論に到達するうえで不要な迂回や重複する補題列を検出する。これにより、カット消去の適用順序や枝刈りの優先順位を動的に調整できる。ただし、カットを実際に消去してよいかどうかは、正規化手続きや証明変換の保存性によって検証されなければならない。

\begin{example}[$\MC$ における冗長性検出]
証明中に同型的な中間命題が複数回現れる場合、LLMはそれらを冗長なカット候補として提示し、正規化手続きが証明の同値性を保ったまま簡約できるかを検証する。
\end{example}

\subsection{$\MT$（識）：統制モナドへの接続}

$\MT$ は、識軸スパイラル $S_{\MT}$ によって五蘊全体を統制するモナドである。ここにLLMを接続すると、LLMは個別の導出や正規化ではなく、五蘊全体のバランスを調整するメタ戦略として機能する。

具体的には、LLMは現在の探索状況に応じて、搾取の検出を強めるべきか（$\Md$）、二重否定の回収を試すべきか（$\Mnn$）、探索戦略を洗練すべきか（$\Mr$）、証明を正規化すべきか（$\MC$）を確率的に判断する。したがって、$\MT$ へのLLM接続は、単一モジュールの補助ではなく、五蘊全体の活性化順序を制御する確率論的統制層を形成する。

\begin{example}[$\MT$ におけるモナド選択]
探索が深くなりすぎて分岐が爆発している場合、LLMは $\Mr$ による分岐選択よりも $\MC$ による枝刈りを優先すべきだと判断する。一方、失敗枝の原因が局所化できない場合には、$\Md$ を活性化して搾取の残余を検出する。
\end{example}

\subsection{統合的な枠組み}

以上をまとめると、モナド五蘊全体にLLMを差し込むことは、各モナドに局所的な確率論的知性を持たせることに相当する。$\Md$ は障害を予測し、$\Mnn$ は否定の回収可能性を探索し、$\Mr$ は分岐を洗練し、$\MC$ は冗長性を検出し、$\MT$ はそれらの活性化順序を統制する。

この構造は、中央集権的な単一LLMではなく、複数の機能的モナドに分散した確率論的知性として理解できる。哲学的には認識の各段階に確率論的要素を導入するモデルであり、認知科学的には思考過程の各フェーズに確率モデルを配置するモデルであり、プログラミング上はアルゴリズムの各モジュールにLLMを組み込む設計原理である。

\section{応用と展望}

\subsection{哲学（認識論）への応用}

エルブラン空間探索層は、認識論のモデルとして応用できる。世界（エルブラン空間）をどのように探索するか（導出戦略・確率論的戦略）は、人間の認識プロセスそのものと見なせる。

\subsection{プログラミング（アルゴリズム設計）への応用}

アルゴリズム設計において、探索空間（エルブラン空間）と探索戦略（導出戦略・確率論的戦略）を明示的に分離することで、より柔軟で効率的なアルゴリズムを設計できる。

\subsection{LLMとの統合}

LLMをエルブラン空間探索層の一部として組み込むことで、論理的厳密性と確率論的柔軟性を両立したAIシステムを構築できる可能性がある。さらに、LLMを $\Mr$ だけでなく $\Md, \Mnn, \MC, \MT$ にも接続すれば、障害予測、二重否定候補の提示、分岐選択、証明正規化、メタ統制を分担する分散型AIシステムとして設計できる。

\section{結論}

本論文では、モナド五蘊の洗練モナド（想）$\Mr$ にエルブラン空間探索層を導入することで、論理的探索（導出・バックトラッキング）と確率論的探索（LLMのような確率モデル）を統一的に扱う枠組みを提案した。さらに、LLMとの接続を想蘊に限定せず、$\Md$ における証明障害予測、$\Mnn$ における二重否定候補の探索、$\MC$ におけるカット消去支援、$\MT$ におけるモナド選択の統制へ拡張する発展的構想を整理した。この枠組みは、顕現論の四句分別およびモナド五蘊の再帰構造と整合的であり、哲学からプログラミングまでを一貫した枠組みで記述できることを示した。

今後の課題として、具体的な圏論的構成（クライスリ圏・EM圏との関係）や、実装レベルのアルゴリズム設計、LLMとの統合実験などが挙げられる。

\begin{thebibliography}{9}

\bibitem{herbrand}
Jacques Herbrand.
\newblock \emph{Recherches sur la th\'eorie de la d\'emonstration}.
\newblock PhD thesis, University of Paris, 1930.

\bibitem{resolution}
J. A. Robinson.
\newblock A Machine-Oriented Logic Based on the Resolution Principle.
\newblock \emph{Journal of the ACM}, 12(1):23--41, 1965.

\bibitem{llm}
Tom B. Brown et al.
\newblock Language Models are Few-Shot Learners.
\newblock \emph{Advances in Neural Information Processing Systems}, 33, 2020.

\bibitem{icat}
千葉将臣.
\newblock 顕現論・一般二諦理論・モナド五蘊に関する基礎論考.
\newblock 未刊行原稿, 2026.

\end{thebibliography}

\end{document}